\documentclass[12pt,a4paper]{article}
\usepackage[utf8]{inputenc}
\usepackage[T2A]{fontenc}
\usepackage[russian]{babel}
\usepackage{amsmath,amssymb}
\usepackage[left=25mm,right=25mm,top=20mm,bottom=20mm]{geometry}
\usepackage{enumitem}

\title{Двухфазная фильтрация нефти с газом в радиальной геометрии.\\
Полная система уравнений с балансом энергии}
\author{}
\date{}

\begin{document}
\maketitle
\tableofcontents
\newpage

%=============================================================================
\section{Постановка задачи}\label{sec:problem}
%=============================================================================

Рассмотрим горизонтальный пласт постоянной мощности~$h$, вскрытый
вертикальной скважиной радиуса~$r_w$. Пласт ограничен контуром
питания~$r = R_{\mathrm{ext}}$, через который поступает жидкость из
удалённой зоны. Задача одномерная --- все величины зависят только от
радиальной координаты~$r$ и времени~$t$:
\begin{equation}\label{eq:domain}
  r \in [r_w,\; R_{\mathrm{ext}}], \qquad t \geqslant 0.
\end{equation}

Пластовый флюид описывается в рамках модели Black~Oil с двумя фазами:
\begin{itemize}[nosep]
  \item \textbf{нефтяная фаза}~(индекс~$o$) --- жидкая смесь тяжёлого
        нефтяного компонента и растворённого в нём лёгкого газового
        компонента;
  \item \textbf{газовая фаза}~(индекс~$g$) --- свободный газ, выделившийся
        из нефти при снижении давления ниже давления насыщения.
\end{itemize}
Водная фаза отсутствует. Испарение нефтяного компонента в газовую фазу
не учитывается: газовая фаза состоит только из газового компонента.

\subsection{Основные допущения}\label{ssec:assumptions}

\begin{enumerate}[label=\arabic*., nosep]
  \item \textbf{Мгновенное локальное фазовое равновесие.}
        При $p < p_b(T)$ газ выделяется из нефти мгновенно; газосодержание
        $R_s(p,T)$ определяется текущими значениями давления и температуры.
  \item \textbf{Единое давление фаз.}
        Капиллярное давление не учитывается: $p_o = p_g \equiv p$.
  \item \textbf{Температурная зависимость PVT.}
        Все свойства фаз ($R_s$, $B_o$, $B_g$, $\rho_o$, $\rho_g$,
        $\mu_o$, $\mu_g$, $c_{p,o}$, $c_{p,g}$) зависят как от давления,
        так и от температуры.
  \item \textbf{Локальное термодинамическое равновесие порода--флюид.}
        Температура породы и флюидов одинакова в каждой точке: $T_s = T_f = T$.
  \item \textbf{Сжимаемость породы.}
        Пористость зависит от давления: $\phi = \phi(p)$.
  \item \textbf{Отсутствие гравитации.}
        Пласт горизонтальный, задача одномерная по~$r$ --- гравитационные
        эффекты не учитываются.
  \item \textbf{Отсутствие теплообмена с кровлей и подошвой.}
        Тепловые потери через границы пласта по вертикали не учитываются.
\end{enumerate}

\subsection{Неизвестные}\label{ssec:unknowns}

Первичные неизвестные задачи --- давление~$p(r,t)$ и температура~$T(r,t)$.
Газонасыщенность~$S_g(r,t)$ определяется алгебраически из уравнений
баланса масс (см.~раздел~\ref{sec:Sg}).


%=============================================================================
\section{Обозначения}\label{sec:notation}
%=============================================================================

Приведём полный список обозначений, используемых в дальнейшем.

\paragraph{Координаты и время.}
\begin{itemize}[nosep]
  \item $r$ --- радиальная координата (расстояние от оси скважины), м;
  \item $t$ --- время, с.
\end{itemize}

\paragraph{Давление и температура.}
\begin{itemize}[nosep]
  \item $p(r,t)$ --- давление (единое для обеих фаз), Па;
  \item $T(r,t)$ --- температура (единая для породы и флюидов), К;
  \item $p_b(T)$ --- давление насыщения (bubble point), Па;
  \item $p_{wf}(t)$ --- забойное давление, Па;
  \item $p_0$ --- начальное пластовое давление, Па;
  \item $T_0$ --- начальная пластовая температура, К.
\end{itemize}

\paragraph{Насыщенности.}
\begin{itemize}[nosep]
  \item $S_o(r,t)$, $S_g(r,t)$ --- объёмные доли нефтяной и газовой фаз
        в поровом пространстве; $S_o + S_g = 1$.
\end{itemize}

\paragraph{PVT-свойства.}
\begin{itemize}[nosep]
  \item $R_s(p,T)$ --- газосодержание нефти (объём растворённого газа при
        стандартных условиях на единицу объёма чистого нефтяного компонента при стандартных условиях),
        м$^3$ст\,/\,м$^3$ст;
  \item $B_o(p,T)$ --- объёмный коэффициент нефти (отношение объёма
        нефтяной фазы при пластовых условиях к объёму нефтяного компонента
        при стандартных), м$^3$пл\,/\,м$^3$ст;
  \item $B_g(p,T)$ --- объёмный коэффициент газа, м$^3$пл\,/\,м$^3$ст;
  \item $\rho_o(p,T)$, $\rho_g(p,T)$ --- плотности фаз при пластовых
        условиях, кг/м$^3$;
  \item $\rho_o^{\text{ст}}$, $\rho_g^{\text{ст}}$ --- плотности
        компонентов при стандартных условиях, кг/м$^3$;
  \item $\mu_o(p,T)$, $\mu_g(p,T)$ --- динамические вязкости фаз, Па$\cdot$с;
  \item $c_{p,o}(p,T)$, $c_{p,g}(p,T)$ --- удельные изобарные теплоёмкости
        фаз, Дж/(кг$\cdot$К);
  \item $\alpha_o(p,T)$, $\alpha_g(p,T)$ --- коэффициенты объёмного
        теплового расширения фаз \textbf{при фиксированном составе},
        К$^{-1}$:
        \begin{equation}\label{eq:thermal_expansion}
          \alpha_\alpha = -\frac{1}{\rho_\alpha}
            \frac{\partial\rho_\alpha}{\partial T}\bigg|_{p,\,R_s},
            \qquad \alpha \in \{o, g\}.
        \end{equation}
        Для нефтяной фазы производная берётся при фиксированном
        газосодержании~$R_s$ (т.\,е.\ без учёта изменения состава при
        изменении температуры). Для газовой фазы (чистый газовый
        компонент) ограничение на состав тривиально;
  \item $L(p,T)$ --- удельная теплота дегазации на единицу массы
        газового компонента, переходящего из растворённого состояния
        в свободное (см.~раздел~\ref{ssec:degas_heat}), Дж/кг газа.
\end{itemize}

\paragraph{Сжимаемости и коэффициенты расширения объёмных факторов.}

Для компактной записи уравнений введём следующие производные объёмных
коэффициентов \textbf{при фиксированном составе фазы}
(фиксированном~$R_s$ для нефтяной фазы):
\begin{equation}\label{eq:co_cg}
  c_o = -\frac{1}{B_o}\frac{\partial B_o}{\partial p}\bigg|_{T,\,R_s},
  \qquad
  c_g = -\frac{1}{B_g}\frac{\partial B_g}{\partial p}\bigg|_T
\end{equation}
--- <<сжимаемости>> нефтяной и газовой фаз (Па$^{-1}$), и
\begin{equation}\label{eq:beta_o_g}
  \beta_o = \frac{1}{B_o}\frac{\partial B_o}{\partial T}\bigg|_{p,\,R_s},
  \qquad
  \beta_g = \frac{1}{B_g}\frac{\partial B_g}{\partial T}\bigg|_p
\end{equation}
--- температурные коэффициенты расширения объёмных факторов (К$^{-1}$).

\paragraph{Замечание.}
Условие $R_s = \mathrm{const}$ принципиально для нефтяной фазы.
В модели Black~Oil корреляция $B_o$ часто задаётся как функция $p$
и~$T$ вдоль кривой равновесия, т.\,е.\ с подставленным $R_s(p,T)$.
В~этом случае полная производная по~$T$ при фиксированном~$p$ включает
вклад от изменения состава:
\begin{equation}\label{eq:beta_o_full}
  \frac{\partial B_o}{\partial T}\bigg|_p
  = \frac{\partial B_o}{\partial T}\bigg|_{p,\,R_s}
    + \frac{\partial B_o}{\partial R_s}\bigg|_{p,\,T}
      \frac{\partial R_s}{\partial T}\bigg|_p.
\end{equation}
В~\textbf{уравнении энергии} используется именно первое слагаемое ---
производная при фиксированном составе ($\beta_o$, $c_o$, $\alpha_o$).
Вклад изменения состава учтён отдельно через теплоту
дегазации~$Q_{\mathrm{deg}}$.

В~\textbf{балансах масс}, напротив, необходимо учитывать полную
производную~$B_o$ вдоль кривой насыщения, что приводит к <<полным>>
коэффициентам~$c_o^{\mathrm{sat}}$, $\beta_o^{\mathrm{sat}}$
(см.~\eqref{eq:co_sat}--\eqref{eq:beta_o_sat}).

\paragraph{Пористая среда.}
\begin{itemize}[nosep]
  \item $\phi(p)$ --- пористость (зависит от давления);
  \item $c_r$ --- сжимаемость порового пространства, Па$^{-1}$:
        \begin{equation}\label{eq:cr}
          c_r = \frac{1}{\phi}\frac{d\phi}{dp};
        \end{equation}
  \item $k$ --- абсолютная проницаемость, м$^2$;
  \item $k_{ro}(S_o)$, $k_{rg}(S_g)$ --- относительные фазовые
        проницаемости;
  \item $\lambda_{\mathrm{eff}}$ --- эффективная теплопроводность
        системы <<порода~+ флюиды>>, Вт/(м$\cdot$К);
  \item $(\rho c_p)_s$ --- объёмная теплоёмкость твёрдого скелета породы,
        Дж/(м$^3\cdot$К).
\end{itemize}

\paragraph{Вспомогательные обозначения.}
\begin{equation}\label{eq:xi_o}
  \xi_o \equiv \frac{S_o}{B_o} = \frac{1 - S_g}{B_o}
\end{equation}
--- объём нефтяного компонента (при ст.~условиях) в единице порового объёма.


%=============================================================================
\section{Закон Дарси}\label{sec:darcy}
%=============================================================================

Скорости фильтрации (скорости Дарси) для каждой фазы:
\begin{equation}\label{eq:darcy}
  u_\alpha = -\frac{k\,k_{r\alpha}}{\mu_\alpha}\,
    \frac{\partial p}{\partial r}, \qquad \alpha \in \{o, g\}.
\end{equation}
Давление убывает от контура питания к скважине, поэтому
$\partial p/\partial r > 0$. Знак <<минус>> в~\eqref{eq:darcy}
обеспечивает $u_\alpha < 0$~--- скорость направлена в сторону
уменьшения~$r$, т.\,е.\ к скважине.

Здесь $k$~--- абсолютная проницаемость, $k_{r\alpha}(S_o)$~--- относительная
фазовая проницаемость фазы~$\alpha$, $\mu_\alpha(p,T)$~--- вязкость.


%=============================================================================
\section{Сжимаемость породы}\label{sec:rock_compr}
%=============================================================================

Пористость зависит от давления. Связь задаётся сжимаемостью порового
пространства~\eqref{eq:cr}:
\begin{equation}\label{eq:dphi_dt}
  \frac{\partial\phi}{\partial t} = \phi\,c_r\,\frac{\partial p}{\partial t},
  \qquad
  \frac{\partial\phi}{\partial r} = \phi\,c_r\,\frac{\partial p}{\partial r}.
\end{equation}
В дальнейшем при раскрытии производных от произведений, содержащих~$\phi$,
всюду используются формулы~\eqref{eq:dphi_dt}.


%=============================================================================
\section{Баланс масс}\label{sec:mass}
%=============================================================================

Выписываем законы сохранения массы для каждого из двух компонентов:
нефтяного и газового. Компонент --- это химическое вещество, которое
может присутствовать в одной или нескольких фазах.

%-----------------------------------------------------------------------------
\subsection{Нефтяной компонент}\label{ssec:mass_oil}
%-----------------------------------------------------------------------------

Нефтяной компонент содержится только в нефтяной фазе. В единице
объёма пористой среды масса нефтяного компонента равна
\[
  \phi\,\frac{S_o}{B_o}\,\rho_o^{\text{ст}}.
\]
Здесь $S_o/B_o = \xi_o$~--- объём нефтяного компонента (при ст.~условиях)
на единицу порового объёма. Множитель $\rho_o^{\text{ст}}$ переводит
объём при ст.~условиях в массу; он постоянен и сокращается. Закон
сохранения в объёмных единицах ст.~условий:
\begin{equation}\label{eq:mass_oil}
  \boxed{
    \frac{\partial}{\partial t}\!\left(\phi\,\frac{S_o}{B_o}\right)
    + \frac{1}{r}\,\frac{\partial}{\partial r}\!
      \left(r\,\frac{u_o}{B_o}\right) = 0.
  }
\end{equation}

\paragraph{Смысл слагаемых.}
Первое слагаемое --- скорость изменения запаса нефтяного компонента
в единице объёма пласта. Второе --- дивергенция потока нефтяного компонента
в радиальной геометрии. Множитель $1/r$ и производная $\partial(r\cdot)/\partial r$
--- стандартная форма оператора дивергенции в цилиндрических координатах
при осевой симметрии.

%-----------------------------------------------------------------------------
\subsection{Газовый компонент}\label{ssec:mass_gas}
%-----------------------------------------------------------------------------

Газовый компонент присутствует в двух фазах:
\begin{itemize}[nosep]
  \item в нефтяной фазе --- в растворённом виде; объём растворённого газа
        (при ст.~условиях) на единицу порового объёма равен $R_s\,\xi_o$;
  \item в газовой фазе --- в свободном виде; объём свободного газа (при
        ст.~условиях) на единицу порового объёма равен $S_g/B_g$.
\end{itemize}

Закон сохранения:
\begin{equation}\label{eq:mass_gas}
  \boxed{
    \frac{\partial}{\partial t}\!\left(\phi\!\left[
      \frac{R_s\,S_o}{B_o} + \frac{S_g}{B_g}
    \right]\right)
    + \frac{1}{r}\,\frac{\partial}{\partial r}\!\left(r\!\left[
      \frac{R_s\,u_o}{B_o} + \frac{u_g}{B_g}
    \right]\right) = 0.
  }
\end{equation}

Первое слагаемое --- скорость изменения полного запаса газового компонента
(растворённого и свободного). Второе --- дивергенция полного потока газового
компонента: растворённый газ переносится вместе с нефтяной фазой со
скоростью~$u_o$, свободный газ --- со скоростью~$u_g$.


%=============================================================================
\section{Раскрытие балансов масс}\label{sec:mass_expand}
%=============================================================================

Для получения уравнений на первичные переменные~$p$ и~$T$ необходимо
раскрыть производные по времени в~\eqref{eq:mass_oil}
и~\eqref{eq:mass_gas}, учитывая зависимости $\phi(p)$, $B_o(p,T)$,
$B_g(p,T)$, $R_s(p,T)$.

%-----------------------------------------------------------------------------
\subsection{Раскрытие баланса нефтяного компонента}\label{ssec:oil_expand}
%-----------------------------------------------------------------------------

Раскроем временную производную в~\eqref{eq:mass_oil}. Используем
правило дифференцирования произведения:
\begin{equation}\label{eq:oil_dt_step1}
  \frac{\partial}{\partial t}\!\left(\phi\,\xi_o\right)
  = \phi\,\frac{\partial\xi_o}{\partial t}
    + \xi_o\,\frac{\partial\phi}{\partial t}.
\end{equation}

Второе слагаемое, с учётом~\eqref{eq:dphi_dt}:
\begin{equation}\label{eq:oil_dt_phi}
  \xi_o\,\frac{\partial\phi}{\partial t}
  = \phi\,c_r\,\xi_o\,\frac{\partial p}{\partial t}.
\end{equation}

Раскроем $\partial\xi_o/\partial t$. Напомним, $\xi_o = (1 - S_g)/B_o$.
По правилу дифференцирования дроби:
\begin{equation}\label{eq:dxi_dt}
  \frac{\partial\xi_o}{\partial t}
  = \frac{1}{B_o}\!\left(-\frac{\partial S_g}{\partial t}\right)
    - \frac{1 - S_g}{B_o^2}\,\frac{\partial B_o}{\partial t}
  = -\frac{1}{B_o}\,\frac{\partial S_g}{\partial t}
    - \frac{\xi_o}{B_o}\,\frac{\partial B_o}{\partial t}.
\end{equation}

Объёмный коэффициент~$B_o$ зависит от $p$, $T$ и состава фазы
(газосодержания~$R_s$). В двухфазной зоне ($p < p_b(T)$) нефть
насыщена: $R_s = R_s(p,T)$, и~$B_o$ определяется как значение на
кривой насыщения: $B_o^{\mathrm{sat}}(p,T) \equiv B_o(p, T, R_s(p,T))$.
Полная временная производная:
\begin{equation}\label{eq:dBo_dt}
  \frac{\partial B_o}{\partial t}
  = \frac{\partial B_o}{\partial p}\bigg|_{T,R_s}\frac{\partial p}{\partial t}
    + \frac{\partial B_o}{\partial T}\bigg|_{p,R_s}\frac{\partial T}{\partial t}
    + \frac{\partial B_o}{\partial R_s}\bigg|_{p,T}
      \frac{\partial R_s}{\partial t}.
\end{equation}

Подставляя определения~\eqref{eq:co_cg}, \eqref{eq:beta_o_g}
и~\eqref{eq:dRs_dt}:
\begin{equation}\label{eq:dBo_dt_compact}
  \frac{1}{B_o}\,\frac{\partial B_o}{\partial t}
  = -c_o\,\frac{\partial p}{\partial t}
    + \beta_o\,\frac{\partial T}{\partial t}
    + \frac{1}{B_o}\frac{\partial B_o}{\partial R_s}\bigg|_{p,T}
      \!\left(
        \frac{\partial R_s}{\partial p}\,\frac{\partial p}{\partial t}
        + \frac{\partial R_s}{\partial T}\,\frac{\partial T}{\partial t}
      \right).
\end{equation}

\paragraph{Замечание: полные и парциальные сжимаемости.}
Последний член в~\eqref{eq:dBo_dt_compact} описывает изменение объёма
нефтяной фазы за счёт изменения её \emph{состава} (растворения или
выделения газа). Удобно ввести <<полные>> коэффициенты вдоль кривой
насыщения, включающие вклад изменения состава:
\begin{align}
  c_o^{\mathrm{sat}} &= c_o
    - \frac{1}{B_o}\frac{\partial B_o}{\partial R_s}\bigg|_{p,T}
      \frac{\partial R_s}{\partial p}\bigg|_T, \label{eq:co_sat} \\
  \beta_o^{\mathrm{sat}} &= \beta_o
    + \frac{1}{B_o}\frac{\partial B_o}{\partial R_s}\bigg|_{p,T}
      \frac{\partial R_s}{\partial T}\bigg|_p. \label{eq:beta_o_sat}
\end{align}
Тогда~\eqref{eq:dBo_dt_compact} принимает компактный вид:
\begin{equation}\label{eq:dBo_dt_sat}
  \frac{1}{B_o}\,\frac{\partial B_o}{\partial t}
  = -c_o^{\mathrm{sat}}\,\frac{\partial p}{\partial t}
    + \beta_o^{\mathrm{sat}}\,\frac{\partial T}{\partial t}.
\end{equation}
Знак в определении~$c_o^{\mathrm{sat}}$ \eqref{eq:co_sat}: при увеличении
давления $R_s$ растёт ($\partial R_s/\partial p > 0$), нефть
поглощает газ и набухает ($\partial B_o/\partial R_s > 0$), поэтому
сжимаемость вдоль кривой насыщения \emph{меньше}, чем при фиксированном
составе: $c_o^{\mathrm{sat}} < c_o$.

Соотношения~\eqref{eq:co_sat}--\eqref{eq:beta_o_sat} связывают
коэффициенты при фиксированном составе (используемые в уравнении
энергии) с коэффициентами вдоль кривой насыщения (появляющимися
в балансе масс). Формула~\eqref{eq:beta_o_full} --- частный случай
\eqref{eq:beta_o_sat} (при $\beta_o^{\mathrm{sat}} = (1/B_o)\,\partial B_o/\partial T|_p$).

Собираем~\eqref{eq:oil_dt_step1}--\eqref{eq:dBo_dt_compact}.
С~учётом~\eqref{eq:dBo_dt_sat} (полная производная вдоль кривой
насыщения):
\begin{equation}\label{eq:mass_oil_expanded}
  \boxed{
    -\frac{\phi}{B_o}\,\frac{\partial S_g}{\partial t}
    + \phi\,\xi_o\,(c_r + c_o^{\mathrm{sat}})\,\frac{\partial p}{\partial t}
    - \phi\,\xi_o\,\beta_o^{\mathrm{sat}}\,\frac{\partial T}{\partial t}
    + \frac{1}{r}\,\frac{\partial}{\partial r}\!
      \left(r\,\frac{u_o}{B_o}\right) = 0.
  }
\end{equation}

\paragraph{Замечание.}
В~\eqref{eq:mass_oil_expanded} фигурируют <<полные>> коэффициенты
$c_o^{\mathrm{sat}}$ и~$\beta_o^{\mathrm{sat}}$
(определённые в~\eqref{eq:co_sat}--\eqref{eq:beta_o_sat}),
поскольку при дифференцировании $B_o$ по времени необходимо учитывать
изменение состава нефтяной фазы вместе с изменением~$p$ и~$T$.
В~уравнении энергии, напротив, используются коэффициенты при
фиксированном составе~$c_o$, $\beta_o$, $\alpha_o$: изменение состава
учтено отдельно через теплоту дегазации~$Q_{\mathrm{deg}}$.

\paragraph{Смысл слагаемых.}
\begin{itemize}[nosep]
  \item Первое: изменение объёма нефтяного компонента из-за изменения
        газонасыщенности (вытеснение нефти газом или наоборот).
  \item Второе: изменение из-за сжатия/расширения при изменении давления.
        Сжимаемость~$c_o^{\mathrm{sat}}$ включает как деформацию фазы
        при фиксированном составе~($c_o$), так и набухание/усадку из-за
        изменения~$R_s$ при изменении~$p$.
  \item Третье: изменение объёма нефтяной фазы при нагреве/охлаждении,
        включая влияние изменения~$R_s(T)$.
  \item Четвёртое: отток нефтяного компонента.
\end{itemize}

%-----------------------------------------------------------------------------
\subsection{Раскрытие баланса газового компонента}\label{ssec:gas_expand}
%-----------------------------------------------------------------------------

Обозначим полный запас газового компонента (в объёмных единицах
ст.~условий на единицу порового объёма):
\begin{equation}\label{eq:xi_g}
  \xi_g = \frac{R_s\,S_o}{B_o} + \frac{S_g}{B_g}
         = R_s\,\xi_o + \frac{S_g}{B_g}.
\end{equation}

Раскроем $\partial\xi_g/\partial t$:
\begin{equation}\label{eq:dxi_g_dt}
  \frac{\partial\xi_g}{\partial t}
  = \xi_o\,\frac{\partial R_s}{\partial t}
    + R_s\,\frac{\partial\xi_o}{\partial t}
    + \frac{1}{B_g}\,\frac{\partial S_g}{\partial t}
    - \frac{S_g}{B_g^2}\,\frac{\partial B_g}{\partial t}.
\end{equation}

Производные $R_s(p,T)$:
\begin{equation}\label{eq:dRs_dt}
  \frac{\partial R_s}{\partial t}
  = \frac{\partial R_s}{\partial p}\,\frac{\partial p}{\partial t}
    + \frac{\partial R_s}{\partial T}\,\frac{\partial T}{\partial t}.
\end{equation}

Производные $B_g(p,T)$ (аналогично~\eqref{eq:dBo_dt_compact}):
\begin{equation}\label{eq:dBg_dt_compact}
  \frac{1}{B_g}\,\frac{\partial B_g}{\partial t}
  = -c_g\,\frac{\partial p}{\partial t}
    + \beta_g\,\frac{\partial T}{\partial t}.
\end{equation}

Подставляя~\eqref{eq:dxi_dt}, \eqref{eq:dRs_dt},
\eqref{eq:dBg_dt_compact} в~\eqref{eq:dxi_g_dt} и группируя
по $\partial S_g/\partial t$, $\partial p/\partial t$,
$\partial T/\partial t$, получаем:
\begin{multline}\label{eq:mass_gas_expanded}
  \boxed{
    \phi\!\left(\frac{1}{B_g} - \frac{R_s}{B_o}\right)
      \frac{\partial S_g}{\partial t}
    + \phi\!\left[\xi_o\,\frac{\partial R_s}{\partial p}
      + \xi_o\,R_s\,(c_r + c_o^{\mathrm{sat}})
      + \frac{S_g}{B_g}\,(c_g + c_r)\right]
      \frac{\partial p}{\partial t}
  } \\
  \boxed{
    + \phi\!\left[\xi_o\,\frac{\partial R_s}{\partial T}
      - R_s\,\xi_o\,\beta_o^{\mathrm{sat}}
      - \frac{S_g}{B_g}\,\beta_g\right]
      \frac{\partial T}{\partial t}
    + \frac{1}{r}\,\frac{\partial}{\partial r}\!\left(r\!\left[
      \frac{R_s\,u_o}{B_o} + \frac{u_g}{B_g}
    \right]\right) = 0.
  }
\end{multline}

\paragraph{Вывод коэффициента при $\partial p/\partial t$.}
Покажем подробно, как получается коэффициент при $\partial p/\partial t$
в~\eqref{eq:mass_gas_expanded}. Вклады дают четыре источника:
\begin{enumerate}[nosep]
  \item из $\xi_o\,\partial R_s/\partial t$: слагаемое
    $\xi_o\,\dfrac{\partial R_s}{\partial p}$;
  \item из $R_s\,\partial\xi_o/\partial t$:
    из~\eqref{eq:dxi_dt} и~\eqref{eq:dBo_dt_sat} при
    фиксированном~$S_g$ следует
    $\partial\xi_o/\partial t\big|_{S_g} =
    \xi_o\,c_o^{\mathrm{sat}}\,\partial p/\partial t
    - \xi_o\,\beta_o^{\mathrm{sat}}\,\partial T/\partial t$;
    вклад по~$\partial p/\partial t$ равен $+R_s\,\xi_o\,c_o^{\mathrm{sat}}$;
  \item из $-(S_g/B_g^2)\,\partial B_g/\partial t$: слагаемое
    $(S_g/B_g)\,c_g$;
  \item из $\xi_g\,\partial\phi/\partial t$: слагаемое $c_r\,\xi_g
    = c_r\,(R_s\,\xi_o + S_g/B_g)$.
\end{enumerate}
Суммируя: $\xi_o\,\partial R_s/\partial p + R_s\,\xi_o\,c_o^{\mathrm{sat}}
+ (S_g/B_g)\,c_g + c_r\,R_s\,\xi_o + c_r\,S_g/B_g
= \xi_o[\partial R_s/\partial p + R_s(c_o^{\mathrm{sat}} + c_r)]
+ (S_g/B_g)(c_g + c_r)$, что даёт коэффициент
в~\eqref{eq:mass_gas_expanded}.


%=============================================================================
\section{Темп выделения свободного газа}\label{sec:degas_rate}
%=============================================================================

Для дальнейшего вывода уравнения энергии нам потребуется выражение для
скорости фазового перехода --- темпа перехода газового компонента из
растворённого состояния в свободное.

Вычтем из уравнения~\eqref{eq:mass_gas}, умноженного на единицу,
произведение~$R_s$ на уравнение~\eqref{eq:mass_oil}:
\begin{multline*}
  \frac{\partial}{\partial t}\!\left(\phi\!\left[
    \frac{R_s\,S_o}{B_o} + \frac{S_g}{B_g}
  \right]\right)
  - R_s\,\frac{\partial}{\partial t}\!\left(\phi\,\frac{S_o}{B_o}\right)
  \\
  + \frac{1}{r}\frac{\partial}{\partial r}\!\left(r\!\left[
    \frac{R_s\,u_o}{B_o} + \frac{u_g}{B_g}
  \right]\right)
  - R_s\,\frac{1}{r}\frac{\partial}{\partial r}\!
    \left(r\,\frac{u_o}{B_o}\right) = 0.
\end{multline*}

Раскроем первые два слагаемых:
\[
  \frac{\partial}{\partial t}\!\left(\phi\,R_s\,\xi_o\right)
  - R_s\,\frac{\partial}{\partial t}\!\left(\phi\,\xi_o\right)
  = \phi\,\xi_o\,\frac{\partial R_s}{\partial t}.
\]
Аналогично для конвективных членов:
\[
  \frac{1}{r}\frac{\partial}{\partial r}\!\left(r\,\frac{R_s\,u_o}{B_o}\right)
  - R_s\,\frac{1}{r}\frac{\partial}{\partial r}\!
    \left(r\,\frac{u_o}{B_o}\right)
  = \frac{u_o}{B_o}\,\frac{\partial R_s}{\partial r}.
\]

В результате получаем уравнение для темпа появления свободного газа:
\begin{equation}\label{eq:free_gas_rate}
  \boxed{
    \frac{\partial}{\partial t}\!\left(\phi\,\frac{S_g}{B_g}\right)
    + \frac{1}{r}\frac{\partial}{\partial r}\!
      \left(r\,\frac{u_g}{B_g}\right)
    = -\phi\,\xi_o\,\frac{\partial R_s}{\partial t}
      - \frac{u_o}{B_o}\,\frac{\partial R_s}{\partial r}.
  }
\end{equation}

Правая часть~\eqref{eq:free_gas_rate} --- это объёмная скорость
появления свободного газа (в м$^3$ст/(м$^3\cdot$с)).
Физический смысл: свободный газ образуется, когда газосодержание~$R_s$
уменьшается, т.\,е.\ когда $\partial R_s/\partial t < 0$ (давление
падает) или $\partial R_s/\partial r > 0$ (нефть втекает из зоны
большего $R_s$ в зону меньшего).

Массовый темп выделения газа:
\begin{equation}\label{eq:mdot}
  \dot{m} = -\rho_g^{\text{ст}}\!\left[
    \phi\,\xi_o\!\left(
      \frac{\partial R_s}{\partial p}\,\frac{\partial p}{\partial t}
      + \frac{\partial R_s}{\partial T}\,\frac{\partial T}{\partial t}
    \right)
    + \frac{u_o}{B_o}\!\left(
      \frac{\partial R_s}{\partial p}\,\frac{\partial p}{\partial r}
      + \frac{\partial R_s}{\partial T}\,\frac{\partial T}{\partial r}
    \right)
  \right].
\end{equation}

При~$\dot{m} > 0$ газ выделяется из нефти; при~$\dot{m} < 0$~---
растворяется обратно.


%=============================================================================
\section{Уравнение давления}\label{sec:pressure}
%=============================================================================

Для получения уравнения на давление~$p$ необходимо исключить
$\partial S_g/\partial t$ из системы~\eqref{eq:mass_oil_expanded},
\eqref{eq:mass_gas_expanded}.

Из~\eqref{eq:mass_oil_expanded} выразим $\partial S_g/\partial t$:
\begin{equation}\label{eq:dSg_dt}
  \frac{\partial S_g}{\partial t}
  = S_o\,(c_r + c_o^{\mathrm{sat}})\,\frac{\partial p}{\partial t}
    - S_o\,\beta_o^{\mathrm{sat}}\,\frac{\partial T}{\partial t}
    + \frac{B_o}{\phi}\,\frac{1}{r}\,\frac{\partial}{\partial r}\!
      \left(r\,\frac{u_o}{B_o}\right).
\end{equation}

Подставим~\eqref{eq:dSg_dt} в~\eqref{eq:mass_gas_expanded}.
После подстановки и группировки (подробности приведены ниже) получаем
уравнение давления:
\begin{equation}\label{eq:pressure}
  \boxed{
    C_p\,\frac{\partial p}{\partial t}
    + C_T\,\frac{\partial T}{\partial t}
    + \mathcal{D} = 0,
  }
\end{equation}
где коэффициенты~$C_p$, $C_T$ и пространственный оператор~$\mathcal{D}$
определены ниже.

\paragraph{Коэффициент при $\partial p/\partial t$:}
\begin{equation}\label{eq:Cp}
  C_p = \phi\!\left[
    S_o\,(c_r + c_o^{\mathrm{sat}})\!\left(\frac{1}{B_g} - \frac{R_s}{B_o}\right)
    + \xi_o\!\left(\frac{\partial R_s}{\partial p}
      + R_s\,(c_r + c_o^{\mathrm{sat}})\right)
    + \frac{S_g}{B_g}\,(c_g + c_r)
  \right].
\end{equation}

\paragraph{Коэффициент при $\partial T/\partial t$:}
\begin{equation}\label{eq:CT}
  C_T = \phi\!\left[
    -S_o\,\beta_o^{\mathrm{sat}}\!\left(\frac{1}{B_g} - \frac{R_s}{B_o}\right)
    + \xi_o\!\left(\frac{\partial R_s}{\partial T}
      - R_s\,\beta_o^{\mathrm{sat}}\right)
    - \frac{S_g}{B_g}\,\beta_g
  \right].
\end{equation}

\paragraph{Пространственный оператор:}
\begin{equation}\label{eq:D_operator}
  \mathcal{D} = \frac{B_o}{B_g}\,\frac{1}{r}\,
    \frac{\partial}{\partial r}\!\left(r\,\frac{u_o}{B_o}\right)
    + \frac{u_o}{B_o}\,\frac{\partial R_s}{\partial r}
    + \frac{1}{r}\,\frac{\partial}{\partial r}\!
      \left(r\,\frac{u_g}{B_g}\right).
\end{equation}

\paragraph{Подробности подстановки.}
При подстановке~\eqref{eq:dSg_dt} в~\eqref{eq:mass_gas_expanded}
член с $\partial S_g/\partial t$ в~\eqref{eq:mass_gas_expanded} равен
\[
  \phi\!\left(\frac{1}{B_g} - \frac{R_s}{B_o}\right)
  \!\left[
    S_o(c_r+c_o^{\mathrm{sat}})\frac{\partial p}{\partial t}
    - S_o\beta_o^{\mathrm{sat}}\frac{\partial T}{\partial t}
    + \frac{B_o}{\phi}\,\frac{1}{r}\frac{\partial}{\partial r}\!
      \left(r\frac{u_o}{B_o}\right)
  \right].
\]
Пространственный вклад от этого члена:
\[
  \left(\frac{B_o}{B_g} - R_s\right)\frac{1}{r}
    \frac{\partial}{\partial r}\!\left(r\frac{u_o}{B_o}\right).
\]
Исходный пространственный член
из~\eqref{eq:mass_gas_expanded} раскладывается как
\[
  R_s\,\frac{1}{r}\frac{\partial}{\partial r}\!\left(r\frac{u_o}{B_o}\right)
  + \frac{u_o}{B_o}\frac{\partial R_s}{\partial r}
  + \frac{1}{r}\frac{\partial}{\partial r}\!\left(r\frac{u_g}{B_g}\right).
\]
Суммируя, члены с $R_s$ сокращаются, и остаётся~\eqref{eq:D_operator}.


%=============================================================================
\section{Уравнение энергии: строгий вывод}\label{sec:energy}
%=============================================================================

%-----------------------------------------------------------------------------
\subsection{Вывод из первого начала термодинамики}\label{ssec:first_law}
%-----------------------------------------------------------------------------

Прежде чем записать уравнение энергии, покажем, почему оно имеет
именно такую структуру: \emph{внутренняя энергия} в члене накопления
и \emph{энтальпия} в конвективном потоке. Эта комбинация --- прямое
следствие первого начала термодинамики для открытой системы.

\paragraph{Первое начало для открытой системы.}
Рассмотрим неподвижный элемент пористой среды объёмом~$\Delta V$,
через который протекает флюид. Первое начало термодинамики для
открытой системы (элемент объёма обменивается массой с окружением)
гласит:
\begin{equation}\label{eq:first_law_open}
  \frac{dE_{\mathrm{total}}}{dt}
  = \dot{Q} - \dot{W}_{\text{внешн}}
    + \sum_{\text{входы}} \dot{m}_{\mathrm{in}}\,
      \widetilde{h}_{\mathrm{in}}
    - \sum_{\text{выходы}} \dot{m}_{\mathrm{out}}\,
      \widetilde{h}_{\mathrm{out}},
\end{equation}
где $E_{\mathrm{total}}$~--- полная внутренняя энергия вещества внутри
элемента, $\dot{Q}$~--- подвод тепла (кондукция), $\dot{W}_{\text{внешн}}$~---
работа системы над окружением (помимо работы проталкивания),
$\widetilde{h}$~--- удельная энтальпия втекающей/вытекающей массы.

Ключевой вопрос: \emph{почему втекающая масса несёт энтальпию~$h$,
а не внутреннюю энергию~$e$?}

\paragraph{Работа проталкивания (flow work).}
Рассмотрим порцию флюида массой $\delta m$ и удельным объёмом
$v = 1/\rho$, которая \emph{втекает} в элемент через границу, где
давление равно~$p_{\mathrm{in}}$. Чтобы эта порция вошла, внешний
флюид должен <<протолкнуть>> её внутрь, совершая работу:
\[
  \delta W_{\mathrm{in}}
  = p_{\mathrm{in}} \cdot \delta V_{\mathrm{in}}
  = p_{\mathrm{in}} \cdot v_{\mathrm{in}} \cdot \delta m.
\]
Эта работа совершается \emph{над системой} (увеличивает её энергию).

Аналогично, когда порция $\delta m$ \emph{вытекает} через границу
с давлением~$p_{\mathrm{out}}$, система совершает работу:
\[
  \delta W_{\mathrm{out}}
  = p_{\mathrm{out}} \cdot v_{\mathrm{out}} \cdot \delta m.
\]

Таким образом, каждая порция массы, входящая в элемент, приносит с
собой свою внутреннюю энергию~$e$ \emph{плюс} окружение совершает
над системой работу проталкивания~$pv$. Суммарный вклад на единицу
массы:
\begin{equation}\label{eq:flow_work}
  e + pv = e + \frac{p}{\rho} = h.
\end{equation}

\textbf{Вывод:} энтальпия в конвективном потоке~--- это не произвольный
выбор, а следствие того, что конвекция неразрывно связана с работой
проталкивания. Внутренняя энергия описывает запас в элементе, энтальпия~---
то, что переносится потоком.

\paragraph{Формализация для одной фазы в пористой среде.}
Рассмотрим фазу~$\alpha$ в единице объёма пористой среды.

Внутренняя энергия фазы в единице объёма пласта:
$\phi\,S_\alpha\,\rho_\alpha\,e_\alpha$.

Конвективный поток внутренней энергии через границу элемента
(на единицу площади в единицу времени): $\rho_\alpha\,e_\alpha\,u_\alpha$,
где $u_\alpha$~--- скорость фильтрации (Дарси).

Работа проталкивания фазы~$\alpha$ (на единицу площади в единицу
времени): $p\,u_\alpha$ (давление~$\times$ объёмный расход фазы через
единицу площади).

Суммарный конвективный перенос энергии фазой~$\alpha$:
\begin{equation}\label{eq:conv_flux_alpha}
  \rho_\alpha\,e_\alpha\,u_\alpha + p\,u_\alpha
  = \rho_\alpha\,u_\alpha\!\left(e_\alpha + \frac{p}{\rho_\alpha}\right)
  = \rho_\alpha\,h_\alpha\,u_\alpha.
\end{equation}

Других форм работы в нашей задаче нет: вязкая диссипация при фильтрации
пренебрежимо мала (числа Рейнольдса $\ll 1$), кинетическая энергия
тоже (скорости фильтрации малы по сравнению со скоростью звука в
флюиде). Подвод тепла --- кондукция по закону Фурье:
$q = -\lambda_{\mathrm{eff}}\,\partial T/\partial r$.

\paragraph{Сборка для многофазной системы.}
Суммируем по фазам и добавляем породу. В радиальной геометрии
закон сохранения энергии принимает вид:
\begin{multline}\label{eq:energy_derivation}
  \frac{\partial}{\partial t}\!\left[
    (1-\phi)\,(\rho c_p)_s\,T
    + \phi\sum_\alpha S_\alpha\,\rho_\alpha\,e_\alpha
  \right]
  + \frac{1}{r}\,\frac{\partial}{\partial r}\!\left(
    r\sum_\alpha \rho_\alpha\,h_\alpha\,u_\alpha
  \right) \\
  = \frac{1}{r}\,\frac{\partial}{\partial r}\!\left(
      r\,\lambda_{\mathrm{eff}}\,\frac{\partial T}{\partial r}
    \right).
\end{multline}

Левая часть: накопление \emph{внутренней} энергии~$+$ дивергенция
потока \emph{энтальпии}. Правая: кондукция. Это и~есть
уравнение~\eqref{eq:energy_cons}, записанное ниже в компактных
обозначениях.

\paragraph{Почему нельзя поставить энтальпию в обе части?}
Формально можно переписать, заменив $e_\alpha = h_\alpha - p/\rho_\alpha$
в~члене накопления и получив $\phi\sum S_\alpha\rho_\alpha h_\alpha$
вместо $\phi\sum S_\alpha\rho_\alpha e_\alpha$. Но тогда в правой
части появится дополнительный член $+\partial(\phi p)/\partial t$~---
это не <<источник тепла>>, а просто перенос из левой части. Физически
форма <<$e$ в накоплении, $h$ в потоке>> --- прямое следствие первого
начала и не содержит таких искусственных членов.

%-----------------------------------------------------------------------------
\subsection{Консервативная форма в компактных обозначениях}%
\label{ssec:energy_conservative}
%-----------------------------------------------------------------------------

Перепишем~\eqref{eq:energy_derivation} в компактном виде. Введём
обозначения для запасённой энергии~$E$ и конвективного потока~$F$:
\begin{equation}\label{eq:energy_cons}
  \frac{\partial E}{\partial t}
  + \frac{1}{r}\,\frac{\partial}{\partial r}\!\left(r\,F\right)
  = \frac{1}{r}\,\frac{\partial}{\partial r}\!\left(
      r\,\lambda_{\mathrm{eff}}\,\frac{\partial T}{\partial r}
    \right).
\end{equation}

\paragraph{Запасённая энергия.}
Внутренняя энергия складывается из энергии твёрдого скелета и энергии
флюидов в порах:
\begin{equation}\label{eq:E_total}
  E = (1 - \phi)\,(\rho c_p)_s\,T
      + \phi\sum_{\alpha\in\{o,g\}} S_\alpha\,\rho_\alpha\,e_\alpha.
\end{equation}
Здесь $e_\alpha = h_\alpha - p/\rho_\alpha$~--- удельная внутренняя
энергия фазы~$\alpha$, $h_\alpha$~--- удельная энтальпия.
Для породы используется приближение: $e_s \approx c_{p,s}\,T$
(несжимаемый скелет, $c_{v,s} \approx c_{p,s}$).

Подставляя $e_\alpha = h_\alpha - p/\rho_\alpha$ и пользуясь
$S_o + S_g = 1$:
\begin{equation}\label{eq:E_expanded}
  E = (1 - \phi)\,(\rho c_p)_s\,T
      + \phi\sum_\alpha S_\alpha\,\rho_\alpha\,h_\alpha
      - \phi\,p.
\end{equation}

\paragraph{Конвективный поток.}
Конвективный перенос энергии осуществляется потоком энтальпии:
\begin{equation}\label{eq:F_total}
  F = \sum_\alpha \rho_\alpha\,h_\alpha\,u_\alpha.
\end{equation}

%-----------------------------------------------------------------------------
\subsection{Энтальпия многокомпонентной нефтяной фазы}\label{ssec:enthalpy_oil}
%-----------------------------------------------------------------------------

Нефтяная фаза --- смесь двух компонентов: тяжёлого нефтяного и лёгкого
газового (растворённого). Её энтальпия зависит не только от $p$ и $T$,
но и от состава, который характеризуется газосодержанием~$R_s$.

Обозначим:
\begin{itemize}[nosep]
  \item $\bar{h}_{oo}(p,T)$ --- парциальная удельная энтальпия
        нефтяного компонента в нефтяной фазе;
  \item $\bar{h}_{go}(p,T)$ --- парциальная удельная энтальпия
        газового компонента в нефтяной фазе.
\end{itemize}

Энтальпия нефтяной фазы в единице объёма пласта:
\begin{equation}\label{eq:oil_enthalpy_components}
  \phi\,S_o\,\rho_o\,h_o
  = \phi\,\frac{S_o}{B_o}\!\left(
      \rho_o^{\text{ст}}\,\bar{h}_{oo}
      + R_s\,\rho_g^{\text{ст}}\,\bar{h}_{go}
    \right).
\end{equation}
Здесь использована связь~\eqref{eq:xi_o} и формула для плотности
нефтяной фазы:
\begin{equation}\label{eq:rho_oil}
  \rho_o = \frac{\rho_o^{\text{ст}} + R_s\,\rho_g^{\text{ст}}}{B_o}.
\end{equation}

Газовая фаза содержит только один компонент (газ), поэтому её энтальпия
просто $h_g(p,T)$.

%-----------------------------------------------------------------------------
\subsection{Перегруппировка по компонентам}\label{ssec:regroup}
%-----------------------------------------------------------------------------

Перегруппируем энтальпию флюидов по компонентам (а не по фазам).
Это позволит явно выделить член фазового перехода.

\paragraph{Энтальпия флюидов в порах:}
\begin{equation}\label{eq:H_fluids}
  \phi\sum_\alpha S_\alpha\rho_\alpha h_\alpha
  = \underbrace{\phi\,\frac{S_o}{B_o}\,\rho_o^{\text{ст}}\,
      \bar{h}_{oo}}_{\text{нефтяной компонент}}
    + \underbrace{\phi\,\rho_g^{\text{ст}}\!\left[
        \frac{R_s\,S_o}{B_o}\,\bar{h}_{go}
        + \frac{S_g}{B_g}\,h_g
      \right]}_{\text{газовый компонент}}.
\end{equation}

\paragraph{Конвективный поток:}
\begin{equation}\label{eq:F_components}
  F = \underbrace{\frac{u_o}{B_o}\,\rho_o^{\text{ст}}\,
        \bar{h}_{oo}}_{\text{нефтяной комп.}}
      + \underbrace{\rho_g^{\text{ст}}\!\left[
          \frac{R_s\,u_o}{B_o}\,\bar{h}_{go}
          + \frac{u_g}{B_g}\,h_g
        \right]}_{\text{газовый комп.}}.
\end{equation}

%-----------------------------------------------------------------------------
\subsection{Выделение теплоты дегазации}\label{ssec:degas_heat}
%-----------------------------------------------------------------------------

Рассмотрим вклад газового компонента в баланс энтальпии. Обозначим
массовые концентрации газового компонента:
\begin{align}
  m_g^{\text{раств}} &= \phi\,\frac{R_s\,S_o}{B_o}\,
    \rho_g^{\text{ст}}, \label{eq:mg_dissolved} \\
  m_g^{\text{своб}} &= \phi\,\frac{S_g}{B_g}\,
    \rho_g^{\text{ст}}. \label{eq:mg_free}
\end{align}

Возьмём временную производную энтальпии газового компонента:
\begin{multline}\label{eq:dHg_dt}
  \frac{\partial}{\partial t}\!\left[
    m_g^{\text{раств}}\,\bar{h}_{go}
    + m_g^{\text{своб}}\,h_g
  \right]
  = \bar{h}_{go}\,\frac{\partial m_g^{\text{раств}}}{\partial t}
    + m_g^{\text{раств}}\,\frac{\partial\bar{h}_{go}}{\partial t}
  \\
    + h_g\,\frac{\partial m_g^{\text{своб}}}{\partial t}
    + m_g^{\text{своб}}\,\frac{\partial h_g}{\partial t}.
\end{multline}

Аналогично, дивергенция конвективного потока газового компонента:
\begin{multline}\label{eq:divFg}
  \frac{1}{r}\frac{\partial}{\partial r}\!\left(r\!\left[
    \frac{R_s\,u_o}{B_o}\,\rho_g^{\text{ст}}\,\bar{h}_{go}
    + \frac{u_g}{B_g}\,\rho_g^{\text{ст}}\,h_g
  \right]\right) \\
  = \bar{h}_{go}\,\frac{1}{r}\frac{\partial}{\partial r}\!
      \left(r\,\frac{R_s\,u_o}{B_o}\,\rho_g^{\text{ст}}\right)
    + \frac{R_s\,u_o}{B_o}\,\rho_g^{\text{ст}}\,
      \frac{\partial\bar{h}_{go}}{\partial r}
  \\
    + h_g\,\frac{1}{r}\frac{\partial}{\partial r}\!
      \left(r\,\frac{u_g}{B_g}\,\rho_g^{\text{ст}}\right)
    + \frac{u_g}{B_g}\,\rho_g^{\text{ст}}\,
      \frac{\partial h_g}{\partial r}.
\end{multline}

Теперь ключевой шаг. Соберём в сумме~\eqref{eq:dHg_dt}
$+$~\eqref{eq:divFg} все члены, содержащие $\bar{h}_{go}$
как множитель:
\begin{equation}\label{eq:hgo_group}
  \bar{h}_{go}\!\left[
    \frac{\partial m_g^{\text{раств}}}{\partial t}
    + \frac{1}{r}\frac{\partial}{\partial r}\!\left(
        r\,\frac{R_s\,u_o}{B_o}\,\rho_g^{\text{ст}}
      \right)
  \right].
\end{equation}
По определению темпа дегазации~\eqref{eq:mdot}, выражение в квадратных
скобках равно~$-\dot{m}$. Следовательно:
\begin{equation}\label{eq:hgo_result}
  \text{вклад с }\bar{h}_{go} = -\bar{h}_{go}\,\dot{m}.
\end{equation}

Аналогично, члены с $h_g$ как множителем:
\begin{equation}\label{eq:hg_group}
  h_g\!\left[
    \frac{\partial m_g^{\text{своб}}}{\partial t}
    + \frac{1}{r}\frac{\partial}{\partial r}\!\left(
        r\,\frac{u_g}{B_g}\,\rho_g^{\text{ст}}
      \right)
  \right].
\end{equation}
По~\eqref{eq:free_gas_rate} (умноженному на $\rho_g^{\text{ст}}$),
выражение в скобках равно~$+\dot{m}$:
\begin{equation}\label{eq:hg_result}
  \text{вклад с }h_g = +h_g\,\dot{m}.
\end{equation}

\paragraph{Сумма~\eqref{eq:hgo_result} и~\eqref{eq:hg_result}:}
\begin{equation}\label{eq:phase_change_heat}
  \boxed{
    (h_g - \bar{h}_{go})\,\dot{m} \equiv L\,\dot{m},
  }
\end{equation}
где введена \textbf{удельная теплота дегазации} (на единицу массы
газового компонента):
\begin{equation}\label{eq:L_def}
  \boxed{
    L(p,T) = h_g(p,T) - \bar{h}_{go}(p,T).
  }
\end{equation}

Здесь:
\begin{itemize}[nosep]
  \item $h_g(p,T)$~--- удельная энтальпия газового компонента в
        свободном состоянии (в газовой фазе), Дж/кг газа;
  \item $\bar{h}_{go}(p,T)$~--- парциальная удельная энтальпия газового
        компонента, растворённого в нефтяной фазе, Дж/кг газа.
        Эта величина зависит от состава раствора, т.\,е.\ от~$R_s$.
        При равновесии $R_s = R_s(p,T)$, поэтому $\bar{h}_{go}$
        выражается как функция~$p$ и~$T$ вдоль кривой равновесия.
\end{itemize}

Физический смысл: $L$~--- энергия, необходимая для перевода единицы
массы газового компонента из растворённого состояния в свободное при
данных $p$ и~$T$. При~$L > 0$ (а это типичный случай, поскольку
свободный газ имеет большую энтальпию, чем растворённый) выделение
газа ($\dot{m} > 0$) требует подвода энергии, т.\,е.\ система
охлаждается.

\paragraph{Замечание о природе~$L$.}
В отличие от теплоты фазового перехода однокомпонентной системы (описываемой
уравнением Клапейрона--Клаузиуса), величина~$L$~--- это разность парциальных
энтальпий компонента в разных фазах многокомпонентной смеси. Уравнение
Клапейрона--Клаузиуса здесь неприменимо. На практике~$L(p,T)$ определяется
из PVT-корреляций или калориметрических измерений при дифференциальном
разгазировании.

%-----------------------------------------------------------------------------
\subsection{Гладкие члены}\label{ssec:smooth_terms}
%-----------------------------------------------------------------------------

После вычленения члена фазового перехода~\eqref{eq:phase_change_heat}
остаются <<гладкие>> члены --- производные энтальпий при фиксированном
составе фаз. Они описывают изменение энтальпии за счёт изменения
$T$ и $p$ без перехода вещества между фазами.

Для нефтяной фазы как целого (при фиксированном~$R_s$) термодинамическое
тождество имеет вид:
\begin{equation}\label{eq:dh_oil}
  dh_o\big|_{R_s} = c_{p,o}\,dT + v_o\,(1 - \alpha_o\,T)\,dp,
\end{equation}
где $v_o = 1/\rho_o$~--- удельный объём нефтяной фазы.

Для газовой фазы:
\begin{equation}\label{eq:dh_gas}
  dh_g = c_{p,g}\,dT + v_g\,(1 - \alpha_g\,T)\,dp.
\end{equation}

Тождества~\eqref{eq:dh_oil} и~\eqref{eq:dh_gas} --- следствие определения
энтальпии $h = h(T,p)$ и соотношения Максвелла. Член $c_{p,\alpha}\,dT$~---
изобарный нагрев; член $v_\alpha(1 - \alpha_\alpha T)\,dp$~---
изотермическое сжатие/расширение.

Гладкие вклады от нефтяной фазы (временная и конвективная части):
\begin{multline}\label{eq:smooth_oil}
  \phi\,S_o\,\rho_o\!\left[
    c_{p,o}\,\frac{\partial T}{\partial t}
    + v_o\,(1-\alpha_o T)\,\frac{\partial p}{\partial t}
  \right] \\
  + \rho_o\,u_o\!\left[
    c_{p,o}\,\frac{\partial T}{\partial r}
    + v_o\,(1-\alpha_o T)\,\frac{\partial p}{\partial r}
  \right].
\end{multline}

Поясним, как получена формула~\eqref{eq:smooth_oil}. Полный
вклад нефтяной фазы в $\partial E/\partial t + \mathrm{div}(rF)/r$
содержит члены вида $h_o\,\partial(\phi S_o\rho_o)/\partial t$
и~$h_o\,\mathrm{div}(\rho_o u_o)/r$. По закону сохранения
массы нефтяной фазы $\partial(\phi S_o\rho_o)/\partial t
+ r^{-1}\partial(r\rho_o u_o)/\partial r = 0$, поэтому эти
слагаемые обращаются в нуль. Остаются материальные производные
энтальпии~$dh_o|_{R_s}$~\eqref{eq:dh_oil}, умноженные на массу
фазы в элементе (для временной части) и массовый расход (для
конвективной). Строго: здесь $h_o$~--- энтальпия фазы как целого,
а~не парциальная энтальпия одного компонента; при $R_s = \mathrm{const}$
состав фазы фиксирован и различие между $h_o$ и~$\bar{h}_{oo}$
корректно обработано через Гиббса--Дюгема.

Поскольку $\rho_o v_o = 1$, члены с давлением
упрощаются:
\begin{equation}\label{eq:smooth_oil_p}
  \phi\,S_o\,(1-\alpha_o T)\,\frac{\partial p}{\partial t}
  + u_o\,(1-\alpha_o T)\,\frac{\partial p}{\partial r}.
\end{equation}

Аналогично для газовой фазы:
\begin{multline}\label{eq:smooth_gas}
  \phi\,S_g\,\rho_g\!\left[
    c_{p,g}\,\frac{\partial T}{\partial t}
    + v_g\,(1-\alpha_g T)\,\frac{\partial p}{\partial t}
  \right] \\
  + \rho_g\,u_g\!\left[
    c_{p,g}\,\frac{\partial T}{\partial r}
    + v_g\,(1-\alpha_g T)\,\frac{\partial p}{\partial r}
  \right],
\end{multline}
с упрощением давленческих членов аналогично~\eqref{eq:smooth_oil_p}.

%-----------------------------------------------------------------------------
\subsection{Вклад породы и члена $-\phi p$}\label{ssec:rock_phip}
%-----------------------------------------------------------------------------

\paragraph{Порода.}
Из~\eqref{eq:E_expanded}:
\begin{multline}\label{eq:rock_dt}
  \frac{\partial}{\partial t}\!\left[
    (1-\phi)\,(\rho c_p)_s\,T
  \right]
  = (1-\phi)\,(\rho c_p)_s\,\frac{\partial T}{\partial t}
  \\
    - (\rho c_p)_s\,T\,\frac{\partial\phi}{\partial t}
  = (1-\phi)\,(\rho c_p)_s\,\frac{\partial T}{\partial t}
    - (\rho c_p)_s\,T\,\phi\,c_r\,\frac{\partial p}{\partial t}.
\end{multline}

\paragraph{Член $-\phi p$.}
Из~\eqref{eq:E_expanded}:
\begin{equation}\label{eq:phip_dt}
  -\frac{\partial(\phi\,p)}{\partial t}
  = -\phi\,\frac{\partial p}{\partial t}
    - p\,\frac{\partial\phi}{\partial t}
  = -\phi\,(1 + p\,c_r)\,\frac{\partial p}{\partial t}.
\end{equation}

%-----------------------------------------------------------------------------
\subsection{Сборка уравнения}\label{ssec:assembly}
%-----------------------------------------------------------------------------

Подставляем все вклады в~\eqref{eq:energy_cons}. Левая часть:
$\partial E/\partial t + \mathrm{div}(rF)/r$.

\paragraph{Члены с $\partial T/\partial t$:}
\begin{equation}\label{eq:T_accum}
  \left[(1-\phi)\,(\rho c_p)_s
    + \phi\,(S_o\,\rho_o\,c_{p,o} + S_g\,\rho_g\,c_{p,g})
  \right]\frac{\partial T}{\partial t}
  \equiv (\rho c_p)_{\mathrm{eff}}\,\frac{\partial T}{\partial t},
\end{equation}
где введена эффективная объёмная теплоёмкость системы:
\begin{equation}\label{eq:rhocp_eff}
  \boxed{
    (\rho c_p)_{\mathrm{eff}}
    = (1-\phi)\,(\rho c_p)_s
      + \phi\,(S_o\,\rho_o\,c_{p,o} + S_g\,\rho_g\,c_{p,g}).
  }
\end{equation}

\paragraph{Члены с $\partial T/\partial r$:}
\begin{equation}\label{eq:T_conv}
  (\rho_o\,c_{p,o}\,u_o + \rho_g\,c_{p,g}\,u_g)\,
    \frac{\partial T}{\partial r}
  \equiv (\rho c_p u)_{\mathrm{conv}}\,\frac{\partial T}{\partial r},
\end{equation}
где
\begin{equation}\label{eq:rhocp_conv}
  \boxed{
    (\rho c_p u)_{\mathrm{conv}}
    = \rho_o\,c_{p,o}\,u_o + \rho_g\,c_{p,g}\,u_g.
  }
\end{equation}

\paragraph{Члены с $\partial p/\partial t$:}
Собираем вклады от породы~\eqref{eq:rock_dt},
члена~$-\phi p$~\eqref{eq:phip_dt} и гладких
членов~\eqref{eq:smooth_oil_p},~\eqref{eq:smooth_gas}:
\begin{multline}\label{eq:p_accum}
  -(\rho c_p)_s\,T\,\phi\,c_r\,\frac{\partial p}{\partial t}
  - \phi(1+p\,c_r)\,\frac{\partial p}{\partial t}
  \\
  + \phi\!\left[S_o\,(1-\alpha_o T)
    + S_g\,(1-\alpha_g T)\right]\frac{\partial p}{\partial t}.
\end{multline}

Используя $S_o + S_g = 1$:
\[
  S_o(1-\alpha_o T) + S_g(1-\alpha_g T) = 1 - T\sum_\alpha S_\alpha\alpha_\alpha.
\]

Подставляя и раскрывая скобки:
\begin{align*}
  &\phi\!\left[1 - T\sum_\alpha S_\alpha\alpha_\alpha
    - 1 - p\,c_r - (\rho c_p)_s\,T\,c_r\right]
    \frac{\partial p}{\partial t} \\
  &= -\phi\!\left[T\sum_\alpha S_\alpha\alpha_\alpha
    + c_r\!\left(p + (\rho c_p)_s\,T\right)\right]
    \frac{\partial p}{\partial t}.
\end{align*}

Этот член переносится в правую часть со сменой знака и становится
источником~$Q_{\mathrm{ad}}$.

\paragraph{Члены с $\partial p/\partial r$:}
Из гладких членов~\eqref{eq:smooth_oil_p},~\eqref{eq:smooth_gas}:
\begin{equation}\label{eq:p_conv}
  \left[u_o\,(1-\alpha_o T)
    + u_g\,(1-\alpha_g T)\right]\frac{\partial p}{\partial r}.
\end{equation}

Этот член переносится в правую часть и становится источником~$Q_{\mathrm{JT}}$.

%-----------------------------------------------------------------------------
\subsection{Итоговое уравнение энергии}\label{ssec:energy_final}
%-----------------------------------------------------------------------------

Переносим все члены, кроме $(\rho c_p)_{\mathrm{eff}}\,\partial T/\partial t$,
в правую часть:
\begin{equation}\label{eq:energy_final}
  \boxed{
    (\rho c_p)_{\mathrm{eff}}\,\frac{\partial T}{\partial t}
    + (\rho c_p u)_{\mathrm{conv}}\,\frac{\partial T}{\partial r}
    = \frac{1}{r}\,\frac{\partial}{\partial r}\!\left(
        r\,\lambda_{\mathrm{eff}}\,\frac{\partial T}{\partial r}
      \right)
    + Q_{\mathrm{JT}} + Q_{\mathrm{ad}} + Q_{\mathrm{deg}},
  }
\end{equation}
где источниковые члены определены ниже.

\paragraph{Эффект Джоуля--Томсона (конвективный).}
При стационарной фильтрации через пористую среду температура изменяется
из-за дросселирования --- необратимого расширения флюида при перепаде
давления:
\begin{equation}\label{eq:QJT}
  \boxed{
    Q_{\mathrm{JT}} = -\sum_\alpha u_\alpha\,(1-\alpha_\alpha T)\,
      \frac{\partial p}{\partial r}.
  }
\end{equation}
Проверим знак. При течении к скважине $u_\alpha < 0$,
$\partial p/\partial r > 0$.

Для \textbf{нефтяной фазы} (жидкость): $\alpha_o T < 1$ при типичных
пластовых условиях ($\alpha_o \sim 10^{-3}$~К$^{-1}$, $T \sim 350$~К,
$\alpha_o T \sim 0{,}3$). Тогда $(1-\alpha_o T) > 0$ и
\[
  -u_o\,(1-\alpha_o T)\,\frac{\partial p}{\partial r}
  = -(-)(\text{+})(\text{+}) > 0.
\]
JT-вклад нефтяной фазы \emph{положителен}: дросселирование
жидкости при~$\alpha T < 1$ приводит к слабому \textbf{нагреву},
а~не охлаждению. Это согласуется со стандартной термодинамикой:
коэффициент Джоуля--Томсона $\mu_{\mathrm{JT}} = v(\alpha T - 1)/c_p < 0$
для жидкости, что означает рост температуры при снижении давления.

Для \textbf{газовой фазы}: у идеального газа $\alpha_g T = 1$ точно
и~вклад в~$Q_{\mathrm{JT}}$ равен нулю. Для реального газа при
пластовых условиях $\alpha_g T$ может быть как больше, так и меньше
единицы в~зависимости от температуры относительно температуры инверсии.

\paragraph{Замечание.}
Наблюдаемое на практике \textbf{охлаждение} при добыче нефти
обусловлено главным образом дегазацией ($Q_{\mathrm{deg}} < 0$),
а~не эффектом Джоуля--Томсона жидкой фазы.

\paragraph{Адиабатический эффект.}
Изменение температуры при локальном изменении давления (при фиксированной
массе флюида в элементе объёма):
\begin{equation}\label{eq:Qad}
  \boxed{
    Q_{\mathrm{ad}} = \phi\!\left[
      T\sum_\alpha S_\alpha\,\alpha_\alpha
      + c_r\!\left(p + (\rho c_p)_s\,T\right)
    \right]\frac{\partial p}{\partial t}.
  }
\end{equation}
Первое слагаемое в скобках --- вклад теплового расширения флюидов:
при~$\partial p/\partial t < 0$ флюиды расширяются адиабатически и
охлаждаются ($Q_{\mathrm{ad}} < 0$). Второе --- вклад сжимаемости
породы, состоящий из двух частей:
\begin{itemize}[nosep]
  \item $p\,c_r$: при снижении давления поры сжимаются, работа деформации
        порового объёма $p\,d\phi$ расходуется на необратимую деформацию
        скелета (уходит из системы), что охлаждает флюид;
  \item $(\rho c_p)_s\,T\,c_r$: при уменьшении $\phi$ объёмная доля
        породы~$(1-\phi)$ растёт, увеличивая эффективную теплоёмкость
        элемента; при фиксированной энтальпии это снижает температуру.
\end{itemize}
Оба вклада от породы обычно малы при типичных
$c_r \sim 10^{-10}$~Па$^{-1}$ ($\sim 10^{-5}$~атм$^{-1}$),
но включены для полноты и корректности размерного анализа.

\paragraph{Теплота дегазации.}
\begin{equation}\label{eq:Qdeg}
  \boxed{
    Q_{\mathrm{deg}} = -L(p,T)\,\dot{m},
  }
\end{equation}
где $L(p,T)$~--- удельная теплота дегазации~\eqref{eq:L_def},
$\dot{m}$~--- массовый темп выделения газа~\eqref{eq:mdot}.

При выделении газа ($\dot{m} > 0$) и $L > 0$: $Q_{\mathrm{deg}} < 0$~---
система охлаждается (энергия расходуется на <<отрыв>> молекул газа от
нефти).

\paragraph{Полное выражение для $\dot{m}$.}
Повторим~\eqref{eq:mdot} для удобства:
\begin{equation}\label{eq:mdot_repeat}
  \dot{m} = -\rho_g^{\text{ст}}\!\left[
    \phi\,\xi_o\!\left(
      \frac{\partial R_s}{\partial p}\,\frac{\partial p}{\partial t}
      + \frac{\partial R_s}{\partial T}\,\frac{\partial T}{\partial t}
    \right)
    + \frac{u_o}{B_o}\!\left(
      \frac{\partial R_s}{\partial p}\,\frac{\partial p}{\partial r}
      + \frac{\partial R_s}{\partial T}\,\frac{\partial T}{\partial r}
    \right)
  \right].
\end{equation}

\paragraph{Важное замечание.}
Член~$Q_{\mathrm{deg}}$ через~\eqref{eq:mdot_repeat} содержит
$\partial T/\partial t$ и $\partial T/\partial r$. Это означает, что
при подстановке~$Q_{\mathrm{deg}}$ в уравнение
энергии~\eqref{eq:energy_final} появляются дополнительные вклады
в эффективную теплоёмкость~$(\rho c_p)_{\mathrm{eff}}$ и конвективный
коэффициент~$(\rho c_p u)_{\mathrm{conv}}$. Перегруппировка этих
членов приведена в разделе~\ref{sec:regroup_deg} (при необходимости).


%=============================================================================
\section{Определение $S_g$ из баланса масс}\label{sec:Sg}
%=============================================================================

Газонасыщенность~$S_g$ определяется из раскрытого баланса нефтяного
компонента~\eqref{eq:dSg_dt}:
\begin{equation}\label{eq:Sg_ode}
  \frac{\partial S_g}{\partial t}
  = S_o\,(c_r + c_o^{\mathrm{sat}})\,\frac{\partial p}{\partial t}
    - S_o\,\beta_o^{\mathrm{sat}}\,\frac{\partial T}{\partial t}
    + \frac{B_o}{\phi}\,\frac{1}{r}\,\frac{\partial}{\partial r}\!
      \left(r\,\frac{u_o}{B_o}\right).
\end{equation}

В однофазной зоне ($p > p_b(T)$) газонасыщенность $S_g = 0$,
состав нефти фиксирован ($R_s = R_s^0 = \mathrm{const}$), и~полные
коэффициенты совпадают с парциальными:
$c_o^{\mathrm{sat}} = c_o$, $\beta_o^{\mathrm{sat}} = \beta_o$.
В двухфазной зоне ($p < p_b(T)$) уравнение~\eqref{eq:Sg_ode} определяет
эволюцию~$S_g$ после того, как $p$ и $T$ найдены из~\eqref{eq:pressure}
и~\eqref{eq:energy_final}.


%=============================================================================
\section{Граничные и начальные условия}\label{sec:bc_ic}
%=============================================================================

%-----------------------------------------------------------------------------
\subsection{Начальные условия}\label{ssec:ic}
%-----------------------------------------------------------------------------

В момент $t = 0$ пласт находится в однородном состоянии:
\begin{equation}\label{eq:ic}
  p(r, 0) = p_0, \qquad T(r, 0) = T_0, \qquad S_g(r, 0) = 0.
\end{equation}
Начальное давление $p_0 > p_b(T_0)$: весь газ растворён в нефти,
свободная газовая фаза отсутствует.

%-----------------------------------------------------------------------------
\subsection{Граничные условия по давлению}\label{ssec:bc_p}
%-----------------------------------------------------------------------------

\paragraph{На забое скважины ($r = r_w$):}
\begin{equation}\label{eq:bc_p_well}
  p(r_w, t) = p_{wf}(t),
\end{equation}
где $p_{wf}(t)$~--- заданная функция времени (забойное давление).

\paragraph{На контуре питания ($r = R_{\mathrm{ext}}$):}
\begin{equation}\label{eq:bc_p_ext}
  p(R_{\mathrm{ext}}, t) = p_0 = \mathrm{const}.
\end{equation}
Через контур питания жидкость поступает в пласт при постоянном давлении.

%-----------------------------------------------------------------------------
\subsection{Граничные условия по температуре}\label{ssec:bc_T}
%-----------------------------------------------------------------------------

\paragraph{На контуре питания ($r = R_{\mathrm{ext}}$):}
\begin{equation}\label{eq:bc_T_ext}
  T(R_{\mathrm{ext}}, t) = T_0.
\end{equation}

\paragraph{На забое скважины ($r = r_w$):}
\begin{equation}\label{eq:bc_T_well}
  \frac{\partial T}{\partial r}\bigg|_{r = r_w} = 0
\end{equation}
--- условие свободного выхода (отсутствие кондуктивного теплового потока).


%=============================================================================
\section{Сводка системы уравнений}\label{sec:summary}
%=============================================================================

Полная замкнутая система состоит из трёх уравнений на три неизвестных
$p(r,t)$, $T(r,t)$, $S_g(r,t)$:

\paragraph{1. Уравнение давления~\eqref{eq:pressure}:}
\[
  C_p\,\frac{\partial p}{\partial t}
  + C_T\,\frac{\partial T}{\partial t}
  + \mathcal{D} = 0.
\]

\paragraph{2. Уравнение энергии~\eqref{eq:energy_final}:}
\[
  (\rho c_p)_{\mathrm{eff}}\,\frac{\partial T}{\partial t}
  + (\rho c_p u)_{\mathrm{conv}}\,\frac{\partial T}{\partial r}
  = \frac{1}{r}\,\frac{\partial}{\partial r}\!\left(
      r\,\lambda_{\mathrm{eff}}\,\frac{\partial T}{\partial r}
    \right)
  + Q_{\mathrm{JT}} + Q_{\mathrm{ad}} + Q_{\mathrm{deg}}.
\]

\paragraph{3. Эволюция газонасыщенности~\eqref{eq:Sg_ode}:}
\[
  \frac{\partial S_g}{\partial t}
  = S_o\,(c_r + c_o^{\mathrm{sat}})\,\frac{\partial p}{\partial t}
    - S_o\,\beta_o^{\mathrm{sat}}\,\frac{\partial T}{\partial t}
    + \frac{B_o}{\phi}\,\frac{1}{r}\,\frac{\partial}{\partial r}\!
      \left(r\,\frac{u_o}{B_o}\right).
\]

\paragraph{Замыкающие соотношения:}
\begin{itemize}[nosep]
  \item Закон Дарси~\eqref{eq:darcy}: $u_\alpha = -(k\,k_{r\alpha}/\mu_\alpha)\,\partial p/\partial r$;
  \item PVT-корреляции: $R_s(p,T)$, $B_o(p,T)$, $B_g(p,T)$, $\rho_\alpha(p,T)$,
        $\mu_\alpha(p,T)$, $c_{p,\alpha}(p,T)$, $\alpha_\alpha(p,T)$;
  \item Теплота дегазации: $L(p,T)$;
  \item Относительные фазовые проницаемости: $k_{ro}(S_o)$, $k_{rg}(S_g)$;
  \item Свойства породы: $\phi(p)$, $k$, $(\rho c_p)_s$, $\lambda_{\mathrm{eff}}$.
\end{itemize}

\paragraph{Граничные условия:}
\begin{itemize}[nosep]
  \item $p(r_w,t) = p_{wf}(t)$, \quad $p(R_{\mathrm{ext}},t) = p_0$;
  \item $T(R_{\mathrm{ext}},t) = T_0$, \quad
        $\partial T/\partial r\big|_{r_w} = 0$.
\end{itemize}

\paragraph{Начальные условия:}
$p(r,0) = p_0$, \quad $T(r,0) = T_0$, \quad $S_g(r,0) = 0$.


%=============================================================================
\section*{Перегруппировка $Q_{\mathrm{deg}}$
  (зарезервировано)}\label{sec:regroup_deg}
%=============================================================================

Раздел зарезервирован для будущей перегруппировки члена
$Q_{\mathrm{deg}}$, содержащего $\partial T/\partial t$
и $\partial T/\partial r$.


\end{document}
